%------------------------------------------------- Chapter 1: context
\chapterseparator{1}{General Context and Environment of the Project}

\chapter{General Context and Environment of the Project}

\section{Introduction}

The rise of e-commerce has made vast product catalogues accessible to consumers worldwide, yet navigating these catalogues remains a friction-filled experience. In specialized retail domains such as organic cosmetics, parapharmacy, and wellness, purchasing decisions depend on nuanced personal criteria (skin type, dermatological concerns, ingredient sensitivity, and application order) that ordinary keyword searches and category menus fail to address adequately.

This chapter situates the project within its industrial and organizational setting. It presents the host enterprise, details the technical and e-commerce environment, outlines the data constraints and catalogue characteristics, formalizes the core engineering problem, and details the temporal organisation and methodology applied during the internship.

\section{Presentation of the host organisation}

\subsection{Company background and core activities}

\companyname\ is a Tunisian digital agency and software enterprise specialising in the architecture, engineering, and delivery of enterprise WordPress applications. The firm covers the complete application lifecycle: requirements engineering, bespoke theme and plugin development, continuous integration, secure cloud deployment, and long-term maintenance of web platforms.

For commercial and transactional platforms, \companyname\ leverages the WooCommerce ecosystem. Rather than an arbitrary constraint, standardizing on WordPress and WooCommerce is a deliberate commercial and technical strategy. WooCommerce powers more than 30\% of global e-commerce deployments, offering small- and medium-sized enterprises in Tunisia and the wider MENA region an agile, extensible, and cost-effective foundation that avoids the prohibitive licensing and operational overhead of proprietary enterprise monoliths.

\subsection{Positioning and service missions}

The operational activity of \companyname\ is structured across three core service axes:
\begin{itemize}
  \item \textbf{End-to-end e-commerce delivery}: Custom storefront development, database schema optimization, payment gateway integration (such as local payment processors), logistics connectors, and production deployment;
  \item \textbf{Custom system extensions}: Development of bespoke WordPress plugins and microservices when off-the-shelf market solutions cannot satisfy custom client workflows;
  \item \textbf{System maintenance and evolution}: Performance profiling, security hardening, high-availability monitoring, and catalogue migrations for existing production systems.
\end{itemize}

The project detailed in this report belongs directly to the second axis. The client required an intelligent, advisory product recommendation capability driven by individual customer concerns. Standard commercial WooCommerce recommendation plugins failed completely because they depend strictly on structured product attributes and tags, which were entirely absent from the client's catalogue.

\subsection{Team organisation, roles, and project workflow}

\companyname\ operates with agile, cross-functional teams where each delivery is managed by dedicated software engineers under the leadership of the company founder. Supervision of this internship was provided directly by the project supervisor, \supervisor, ensuring strategic alignment with client specifications.

The project team comprised two engineers: the lead web engineer responsible for storefront deployment, infrastructure, and production database exports; and the author of this report, responsible for the end-to-end recommendation system, from exploratory catalogue audit to agentic LLM orchestration. Table~\ref{tab:roles} details the division of responsibilities, and Figure~\ref{fig:company_workflow} illustrates the interaction workflow between the host company, client infrastructure, and the development pipeline.

\begin{table}[H]
\caption{Division of responsibilities on the project}
\label{tab:roles}
\centering
\begin{tabularx}{\textwidth}{L{4.4cm} X}
\toprule
\textbf{Role} & \textbf{Responsibility} \\
\midrule
Project supervisor (\supervisor) & Strategic framing, client relation, technical governance, and final review \\
Web developer (colleague) & WordPress storefront maintenance, WooCommerce hosting environment, production database exports \\
Intern / Author & Catalogue audit, product data model design, automated extraction pipeline, deterministic scoring engine, FastAPI microservice, WordPress plugin, conversational agent, validation protocol \\
\bottomrule
\end{tabularx}
\end{table}

\placeholderscreenshot{company_workflow.png}{Host organisation and project operational workflow}{fig:company_workflow}

\section{General and technological context}

\subsection{E-commerce in Tunisia and the beauty and wellness sector}

The Tunisian e-commerce market has expanded substantially over recent years, propelled by increasing digital adoption, local parcel delivery networks, and consumer interest in health, parapharmaceutical, and organic products. The client business operates in this thriving segment via \clientsite, an established Tunisian online shop marketing certified natural cosmetics, organic skincare, food supplements, and eco-friendly wellness goods.

The retail catalogue of \clientsite\ reflects local market specificities: it carries 41 distinct regional and international cosmetic and wellness brands. Because no single manufacturer dominates the store, recommendations must operate objectively across multiple competing brands rather than promoting a single cosmetic product line.

\subsection{The WooCommerce application architecture}

The client's technical foundation is WordPress with the WooCommerce extension, running on MariaDB. In this architecture, all catalogued items are stored in the WordPress \texttt{posts} table under the \texttt{product} post-type, with custom attributes serialized into the relational \texttt{postmeta} table and taxonomic categorizations partitioned across \texttt{terms}, \texttt{term\_taxonomy}, and \texttt{term\_relationships}.

While this relational schema supports flexible inventory display and checkout, it imposes significant computational barriers for multi-attribute semantic querying, particularly when meta attributes are sparse or unstructured.

\section{Data constraints and catalogue selection}

\subsection{Catalogue structure and available dataset}

The dataset made available for this project consists of an anonymized SQL dump from the production MariaDB database of \clientsite. The complete export contains:
\begin{itemize}
  \item \textbf{1{,}073 total product entries}, comprising published goods, archived drafts, private tests, and composite pack listings;
  \item \textbf{612 active, published products} eligible for commercial recommendation;
  \item \textbf{69 product categories}, heavily skewed toward anatomical formatting (e.g. \emph{Visage}, \emph{Corps}, \emph{Cheveux}) and merchandising banners (e.g. \emph{Accueil / Homepage} with 239 redundant items);
  \item \textbf{52 commercial packs} marketed as bundled beauty routines;
  \item \textbf{Zero product tags} registered in the database;
  \item \textbf{Skin-type attributes populated on exactly 6 products out of 612}.
\end{itemize}

\subsection{Technical and operational constraints}

Three non-negotiable operational constraints governed the architecture and engineering decisions:
\begin{enumerate}
  \item \textbf{Strict read-only access to the production catalogue}: The recommendation system must never alter or lock production database tables. All analytical exploration and pipeline extraction was performed on a local mirror, and production integration was restricted to read-only queries;
  \item \textbf{Zero-budget API constraint for inference}: The project operated without a funded commercial LLM subscription. This required designing a deterministic rule-based core that operates at zero marginal cost, reserving LLM usage for high-level semantic interpretation and phrasing on free-tier APIs;
  \item \textbf{Absence of an on-site dermatological expert}: No professional cosmetologist was available on-site for immediate real-time validation. Consequently, a systematic 22-profile empirical validation framework was designed to assess recommendation safety and coherence objectively.
\end{enumerate}

\section{Problem statement}

Online shoppers visiting a specialized beauty boutique typically express problems in symptomatic, natural vocabulary (for example, \emph{``My skin feels tight after washing, but I break out around my forehead, and I have dark sun spots''}). Conversely, conventional e-commerce systems require the customer to manually select category filters, deduce suitable active ingredients, and independently construct a coherent application routine (cleanser, toner, serum, moisturizer, sunscreen).

This creates a severe semantic gap:
\begin{itemize}
  \item \textbf{The vocabulary mismatch}: Customers know their symptoms; catalogues list technical chemical formulations and biological mechanisms;
  \item \textbf{The attribute sparsity barrier}: Standard recommender algorithms (collaborative filtering or attribute-based recommenders) fail completely because interaction logs are sparse and structured metadata is absent;
  \item \textbf{The routine coherence requirement}: Unlike standalone book or electronics recommendations, cosmetics require an ordered sequence of complementary products where chemical incompatibilities and contraindications must be respected.
\end{itemize}

The central engineering challenge of this project was to design, implement, and validate an automated system capable of transforming unstructured commercial French text into a structured, queryable knowledge base, and orchestrating a conversational agent grounded in this data to advise shoppers safely.

\section{Project temporal organisation and methodology}

\subsection{Working method and short feedback loops}

Rather than following a rigid, single-pass waterfall model, the project adopted an empirical, short-cycle iterative methodology aligned with CRISP-DM principles (detailed in Section~2.6). Each technical sprint produced an operational, testable prototype. Inspecting the prototype against realistic customer queries immediately exposed hidden data anomalies, which fed directly into the next development cycle.

Figure~\ref{fig:method} illustrates this empirical feedback loop: early UI prototypes revealed critical edge cases (such as lip scrubs erroneously suggested as facial cleansers, or dietary supplements mixed indiscriminately with topical night creams) that purely theoretical database inspections had failed to uncover.

\begin{figure}[H]
\centering
\begin{minipage}{0.92\textwidth}
\begin{lstlisting}[frame=single, basicstyle=\ttfamily\footnotesize]
   Audit & Extract Data  ───►  Build Runnable Prototype  ───►  Inspect UI Output
           ▲                                                         │
           │                                                         ▼
           └──────────  Refine Model & Constraints  ◄──────  Detect Behavioral Flaw
\end{lstlisting}
\end{minipage}
\caption{Empirical iterative development cycle applied throughout the project}
\label{fig:method}
\end{figure}

\subsection{Chronology of project tasks}

The internship spanned five weeks of intensive research, design, coding, and validation. Table~\ref{tab:chronology} outlines the key technical milestones achieved across each project stage, and Figure~\ref{fig:gantt_schedule} presents the temporal Gantt schedule of the project.

\begin{table}[H]
\caption{Chronology of tasks performed across the internship period}
\label{tab:chronology}
\centering
\begin{tabularx}{\textwidth}{L{2.4cm} X}
\toprule
\textbf{Stage} & \textbf{Key technical deliverables and milestones} \\
\midrule
Week 1 & Restoration of the 391\,MB MariaDB database dump; exploratory SQL auditing of schema, taxonomies, and meta tables; quantitative evaluation of attribute sparsity; authoring of the comprehensive Catalogue Audit Report. \\
Week 2 & Conception of the multi-dimensional product data model; implementation of the regex-based extraction pipeline (\texttt{2\_extract.py}); TF-IDF algorithmic resolution of pack contents (\texttt{3\_resolve\_packs.py}); extraction quality assurance reporting. \\
Week 3 & Implementation of the deterministic scoring algorithm and 4-step routine assembly engine (\texttt{scoring.py}); construction of the interactive web demo prototype (\texttt{demo/index.html}); defect diagnosis and pattern refinement. \\
Week 4 & Development of the asynchronous FastAPI REST service (\texttt{service/}); live pricing query integration; implementation of the custom WordPress / WooCommerce shortcode plugin; initial NLU query parsing. \\
Week 4--5 & Design and implementation of the multi-turn conversational agent with ReAct tool-calling loops; safety guardrails and anti-hallucination isolation; 22-profile empirical validation protocol; documentation and LaTeX report redaction. \\
\bottomrule
\end{tabularx}
\end{table}

\placeholderscreenshot{gantt_schedule.png}{Project temporal organisation and timeline of the internship stages}{fig:gantt_schedule}

\section{Conclusion}

This chapter detailed the organizational and industrial setting of the internship at \companyname, working on the production catalogue of \clientsite. The project operated under three fundamental constraints: read-only database security, zero recurring inference budget, and the absence of a resident domain expert. The subsequent chapters demonstrate how these concrete real-world constraints directly dictated the theoretical state-of-the-art positioning, the analytical framework, and the layered implementation architecture of the final system.

\cleardoublepage
